\chapter{Specification fonctionnelle detaillee}
\label{ch:features}

\section{Module Accueil}

La vue \texttt{welcome.blade.php} assure le \textbf{positionnement produit}, les acces rapides (hotels, restaurants, carte) et integre \texttt{<x-chatbot />}. Objectif : conversion vers exploration ou inscription.

\screenshot{figures/screenshots/welcome.png}{Page d'accueil}{Figure -- Landing page Sultan Morocco}

\section{Module Hotels}

\subsection{Regles metier}

\begin{itemize}
  \item Catalogue de \textbf{30} etablissements references au Maroc (Marrakech, Agadir, Casablanca, etc.).
  \item Tri \texttt{best\_value} : note decroissante puis volume d'avis.
  \item Tri \texttt{popularity} : volume d'avis decroissant.
  \item Pagination : \texttt{offset}, \texttt{limit} $\in [1,100]$.
  \item Filtre ville cote client derive de l'URL TripAdvisor embarquee.
\end{itemize}

\subsection{Structure de donnee (extrait)}

\begin{lstlisting}[caption={Champs type d'un hotel (HotelController)}]
"name", "key", "accommodation_type", "url",
"review_summary" => ["rating", "count"],
"price_ranges" => ["minimum", "maximum"],
"geo" => ["latitude", "longitude"],
"image", "merchandising_labels"
\end{lstlisting}

\animHotelLoad

\screenshot{figures/screenshots/hotels.png}{Grille hotels}{Figure -- Listing paginé avec filtres}

\subsection{Enrichissement adresse (Nominatim)}

File d'attente JS ($\sim$1,15 s entre appels) respectant la politique OSM. Coordonnees issues du JSON statique ; adresse affichee est \textbf{calculee a la volee}.

\section{Module Restaurants}

\textbf{15} etablissements, champs : \texttt{city}, \texttt{cuisine}, \texttt{price\_level}, \texttt{specialties}, \texttt{tripadvisor\_url}, image Unsplash.

API \texttt{GET /api/restaurants} : filtre \texttt{city}, tri \texttt{rating|price\_asc|price\_desc}.

\screenshot{figures/screenshots/restaurants.png}{Restaurants}{Figure -- Page restaurants}

\section{Module Carte}

\begin{itemize}
  \item Tuiles : OpenStreetMap + CARTO (attribution requise).
  \item Marqueurs : coordonnees \texttt{geo} des hotels.
  \item Clustering : MarkerCluster (CDN).
  \item Sidebar : liste synchronisee, popup au clic.
\end{itemize}

\animMapFlow
\screenshot{figures/screenshots/map.png}{Carte}{Figure -- Carte Leaflet}

\section{Module Chatbot}

Intents : \texttt{greeting}, \texttt{help}, \texttt{restaurants}, \texttt{hotels}, \texttt{surf}, \texttt{itinerary}, \texttt{weather}. Reponses : \texttt{text}, \texttt{buttons}, \texttt{cards}. Pas de persistance DB.

\animChatbot
\screenshot{figures/screenshots/chatbot.png}{Chatbot}{Figure -- Widget assistant}

\section{Module Authentification}

Signup : nom, email unique, mot de passe $\geq$ 8 caracteres, confirmation. Login : session regeneree. Logout : invalidation complete.

\animAuthFlow
\screenshot{figures/screenshots/login.png}{Login}{Figure -- Authentification}
\screenshot{figures/screenshots/signup.png}{Inscription}{Figure -- Creation de compte utilisateur}

\section{Module Profil et messagerie}

Profil public par ID. Messagerie : thread 1:1, preview dernier message, compteur non-lus, \texttt{read\_at} a l'ouverture.

\screenshot{figures/screenshots/profile.png}{Profil utilisateur}{Figure -- Fiche profil avec acces a la messagerie}

\animMessaging
\screenshot{figures/screenshots/messages-inbox.png}{Inbox}{Figure -- Boite de reception}
\screenshot{figures/screenshots/messages-thread.png}{Thread}{Figure -- Conversation}

\section{Exigences non fonctionnelles}

\begin{table}[H]
\centering
\caption{Exigences non fonctionnelles principales}
\begin{tabular}{@{}lll@{}}
\toprule
\textbf{ID} & \textbf{Exigence} & \textbf{Cible} \\
\midrule
NFR-01 & Temps chargement listing & $<$ 500 ms API locale \\
NFR-02 & Disponibilite listings & 100\% (statique) \\
NFR-03 & Responsive UI & Mobile + desktop \\
NFR-04 & Accessibilite couleurs & Theme clair/sombre \\
NFR-05 & Securite auth & Session + CSRF formulaires \\
\bottomrule
\end{tabular}
\end{table}
